% ============================================================
%  Section 6: Conclusion
% ============================================================

\section{Conclusion}
\label{sec:conclusion}

We presented \agentforge{}, a multi-agent framework for
autonomous software engineering that decomposes code
generation, testing, debugging, and review into five
specialized agents connected by a structured orchestration
pipeline. By executing generated code in a genuine
Docker sandbox and grounding each agent call in retrieved
repository context and episodic memory, \agentforge{}
closes the gap between single-shot LLM code generation
and the iterative, feedback-driven process of human
software engineering.

On \swebench{} Lite, \agentforge{} achieves \RESULT{XX\%}
task resolution, outperforming single-agent and ReAct
baselines by \RESULT{XX} percentage points. Our ablation
study shows that each agent contributes meaningfully,
with the Tester--Debugger loop providing the largest
individual gain — evidence that verified execution
feedback, rather than model scale alone, is the critical
ingredient for strong performance on real-world
software engineering tasks.

\paragraph{Limitations.}
\agentforge{} inherits several limitations of current
LLM-based systems. First, it operates at the file level:
tasks requiring coordinated changes across many files or
build system modifications remain challenging.
Second, the evaluation oracle is binary — a task either
resolves or does not — and does not capture partial credit
for patches that fix the reported symptom but introduce
latent regressions.
Third, all experiments use GPT-4o; performance with
open-weight models is an important open question.

\paragraph{Future Work.}
Promising directions include: (1) multi-file atomic patch
application with rollback on failure; (2) function-level
rather than file-level repository indexing for finer
retrieval granularity; (3) a persistent cross-session
memory store that accumulates repository-specific
knowledge over time; (4) evaluation on additional
benchmarks including HumanEval-Pro and
multi-language SWE-bench variants; and
(5) fine-tuning specialized smaller models for each
agent role to reduce inference cost.

\paragraph{Broader Impact.}
Autonomous software engineering agents have the potential
to substantially accelerate software development,
making programming assistance more accessible to
non-expert developers. At the same time, they introduce
risks around the generation of insecure or incorrect
code at scale. The sandboxed execution environment in
\agentforge{} mitigates the latter risk in the
evaluation setting, but production deployment requires
additional safeguards including human oversight,
static analysis integration, and audit logging.
